\documentclass[a4paper,cleardoubleempty,BCOR1cm]{scrbook}
% use to waste space:
% \documentclass[12pt,a4paper]{article}

% if you have this style and like it.
%\documentclass{acmsiggraph}
%\documentclass[review]{acmsiggraph}      % review
%\documentclass[widereview]{acmsiggraph}  % wide-spaced review
%\documentclass[preprint]{acmsiggraph}    % preprint

% define a \comment{this is a comment which can have linebreaks in it}
\newcommand{\comment}[1]{}
% \newcommand{\todo}[1]{\marginpar{\bf{#1}}}
\newcommand{\todo}[1]{{\color{red}\bf{TODO: #1}}}

\usepackage{mathptmx}
\usepackage[pdftex]{graphicx}
\usepackage[pdftex]{color}
\definecolor{rot}{RGB}{165,30,55} %rote Farbe
\graphicspath{{./images/}}
\usepackage{parskip}
\usepackage{amsmath}
\usepackage{dsfont}
\usepackage{pxfonts}

\usepackage[T1]{fontenc}
\usepackage{textcomp}

% comment these two lines out if you don't want minion/myriad fonts.
% \usepackage[minionint,mathlf]{MinionPro}
% \renewcommand{\sfdefault}{Myriad-LF}
%\usepackage{Myriad}

% no page number on float pages, fixes problems with overlarge diagrams.
\usepackage{fancyhdr}
\pagestyle{fancy}
%\lhead{}
%\chead{}
%\rhead{}
%\lfoot{}
\fancyhf{}
\fancyhead[EL]{\nouppercase{\leftmark}}
\fancyhead[OR]{\nouppercase{\rightmark}}
\cfoot{}
%\fancyfoot[EL]{\iffloatpage{}{\thepage}}
%\fancyfoot[OR]{\iffloatpage{}{\thepage}}
\fancyfoot[EL]{\thepage}
\fancyfoot[OR]{\thepage}
\renewcommand{\headrulewidth}{0pt}
\renewcommand{\footrulewidth}{0pt}

%\usepackage{natbib}		% textual referencing
%\usepackage[numbers,super]{natbib}	% nice superscripts
%\bibliographystyle{chicago}	% shitty
\bibliographystyle{alpha}	% abbr names and year in \cite
%\bibliographystyle{agsm}	% australian, need natbib
%\bibliographystyle{kluwer}	% need natbib
%\bibliographystyle{apalike}	% lengthly
%\bibliographystyle{abbrv}	% minimal?

% use for german line breaking:
%\usepackage[ngerman]{babel}
\usepackage[T1]{fontenc}
\usepackage[utf8x]{inputenc}

% avoid us-style text color destruction:
\frenchspacing
\usepackage{microtype}

% have a nice framebox with border directly around the image:
\fboxsep 0pt
\newcommand{\fimg}[2]{\fbox{\includegraphics[width=#1]{#2}}}

\usepackage{theorem}
\theorembodyfont{\upshape}
\newtheorem{definition}{Definition}

\usepackage{listings}
\lstset{numbers=left, numberstyle=\tiny, basicstyle=\tiny, language=C++}
\usepackage[boxruled]{algorithm2e}
%\usepackage{hyperref}
\usepackage{url}
\usepackage{subfig}

\def\code#1{{\tt{#1}}}



\title{Thesis Template}
\author{Denis Heid \thanks{e-mail: denis.heid@student.uni-tuebingen.de}}
\date{\today}
\begin{document}

\begin{tabular}{lr}
% \includegraphics[width=0.5\linewidth]{logo_sw} % logo bw
 \includegraphics[width=0.5\linewidth]{UT_WBMW_Rot_4C} % logo red
 & \hspace{0.2\linewidth}
 \parbox{0.5\linewidth}{
   \large\bf\textsf{\color{rot}{Mathematisch-\\Naturwissenschaftliche\\Fakultät\\\\}}
   \hspace{-.144cm}\normalsize\textsf{\color{rot}{Computergrafik}}
   \vspace{0.6cm}
 }
\end{tabular}

\vspace*{10ex}

\includegraphics[width=1.055\linewidth]{teaser.png}

Masterarbeit

{\huge\bf\textsf{points2mesh:\\ Learning surface reconstruction from point cloud data}}

%\vspace*{30ex}
\vspace*{14ex}

Eberhard Karls Universität Tübingen\\
Mathematisch-Naturwissenschaftliche Fakultät\\
Wilhelm-Schickard-Institut für Informatik\\
Computergrafik\\
Denis Jan Heid,~ \verb+denis.heid@student.uni-tuebingen.de+,~ 2019

\vspace*{5ex}

\begin{tabular}{@{}l@{\hspace{2em}}l}
  Bearbeitungszeitraum:& Januar 2019-Juli 2019\vspace*{5ex} \\
  Betreuer/Gutachter:& Prof. Dr. Hendrik Lensch, Universität Tübingen\\
  Zweitgutachter:& Prof. Dr. Andreas Schilling, Universität Tübingen
\end{tabular}

\thispagestyle{empty}
\newpage

\chapter*{Selbstst\"andigkeitserkl\"arung}
Hiermit versichere ich, dass ich die vorliegende Masterarbeit selbst\"andig und
nur mit den angegebenen Hilfsmitteln angefertigt habe und dass alle Stellen,
die dem Wortlaut oder dem Sinne nach anderen Werken entnommen sind,
durch Angaben von Quellen als Entlehnung kenntlich gemacht worden sind.
Diese Masterarbeit wurde in gleicher oder \"ahnlicher Form in keinem anderen
Studiengang als Pr\"ufungsleistung vorgelegt.

\vspace*{8ex}
\hrule
\vspace*{2ex}
Denis Heid (Matrikelnummer 3827662), \today



\chapter*{Abstract}
    This thesis introduces \emph{points2mesh}, a novel approach to transform point cloud data into fully-fledged, watertight meshes. 
    Adapted from \emph{pixel2mesh} \cite{wang2018pixel2mesh}, transforming images to meshes, this neural network learns features from unstructured points
    in three-dimensional space and deforms a basic 3D ellipsoidal shape into the final watertight mesh based on these learned features.

    \emph{points2mesh} offers several distinct advantages. As the first deep learning-based method to transform point
    cloud data to watertight meshes while producing competitive results with low inference time, it is a pioneer in its subject-matter.
    Additionally, a \emph{supervised} as well as an \emph{unsupervised} variant is presented.


\chapter*{Zusammenfassung}
    Diese Arbeit stellt \emph{points2mesh} vor, einen neuartigen Ansatz zur Transformation von Punktwolkendaten in manngifaltige Polygonnetze. 
    Dieses neuronale Netzwerk, adaptiert von \emph{pixel2mesh} \cite{wang2018pixel2mesh}, welches Bilder in Polygonnetze umwandelt, lernt Merkmale von unstrukturierten Punkten 
    im dreidimensionalen Raum und verformt ein einfaches 3D Polygonnetz in das endgültige mannigfaltige Polygonnetz, basierend auf diesen erlernten Merkmalen.

    \emph{points2mesh} bietet mehrere Vorteile. Als erste lernbasierte Methode zur Transformation von Punktwolkendaten in mannigfaltige Polygonnetze
    bei gleichzeitiger Erzielung wettbewerbsfähiger Ergebnisse mit geringer Inferenzzeit ist sie ein Pionier in ihrem Fachgebiet.
    Zusätzlich wird sowohl eine beaufsichtigte als auch eine unbeaufsichtigte Lernvariante vorgestellt.


\chapter*{Acknowledgments}
    First and foremost, I would like to thank Fabian Groh for guiding me through this thesis, lending valuable feedback and providing 
    a stream of possible improvements or new avenues to consider. Also, I would like to thank Professor Lensch for allowing me to write
     the thesis at his chair of computer graphics.
    I also have enormous gratitude for my family and Friederike Hanssen, supporting and encouraging me throughout this period.



\tableofcontents

%% braucht kein Mensch ...
%% aber ich mag die :(
%\listoffigures
%\listoftables


% write content here or...
\chapter{Introduction}


\begin{enumerate}
    \item robotics, interact with real world
    \item understand env
    \item what is it, they are working with
    \item autonomous driving 
    \item thus object analysis 
    \item digitize objects somehow
    \item round about of different kinds already present
    \item List different representations of objects (meshes, pointcloud, voxels,structured light, kinect depth..)
    \item lots of data readily available
    \item present novel network configurations deep learning based
    \item end to end solution for shape inference from PC to mesh
    \item some modified loss functions for good results
    \item determine how well it is doing this functions
    \item compare to work doing similar things
    \item non learning based/ learning based
    \item maybe small user study??
  \end{enumerate}
\section{Problem Statement}
mathmatic formulation
\section{chapter recapitulation}

x`\chapter{Background}
\label{sec:background}
Often in computer graphics, it is necessary to process three-dimensional objects. There are many ways to acquire data of the surface of such geometry and more so many techniques to transform and augment that data to a distinct representation. This often non-trivial task is crucial to further process the object in question in later stages of their respective pipelines. Over the years many representations of acquisition data formats and transformation methods, as well as target data formats, have accumulated.\\
In this chapter, various of these techniques and data formats are examined, of which some of them are used in this work as a vehicle for a novel data transformation routine.\\
Initially, classical approaches are examined in section \ref{classic_approaches} which do not rely on artificial intelligence or machine learning methods.
With the recent advances in deep learning, many new approaches have been developed and thus considered in this work. Hence, the general concept of deep learning and neural networks are then reviewed in section \ref{ml_review}.
Subsequently, these more specific machine learning based methods are reassessed in section \ref{ml_approaches} which rely on distinctive statistical features in their initial data format or the dataset itself, thus allowing for the transformation.
\todo{citations for this stuff?}
\section{Classical approaches}
\label{classic_approaches}
\todo{different name}
In this chapter, classical approaches are examined which do not rely on machine
 learning to transform data of objects from one format to another. 

In general, a three-dimensional object may be specified by its representation of the
 surface. This representation varies from unstructured point cloud data in three-dimensional space
  to a graph-based representation like triangle-/ or quad-meshes and even poly meshes
   with more than four neighbors per node. Surface information may also be described
    with voxelized surface representations.  Often the information on an object's
     surface is only partially defined, where some information may be missing. 
     This can be due to acquisition methods, where parts of the object are partially 
     concealed or recorded from insufficient many sides.



\subsection{Transformation methods}
\todo{maybe not needed. only classic approaches as section?}
Starting from these representations of an object's surface, there are many methods 
of transformation to reach another state of representation. 
Transforming the data not only makes further processing easier
 but rather during their process, more data on the surface is 
 computed based on the given input data.

 This task is non-trivial since, from one representation of an object's surface, many reconstructions in the form of another representation are possible. The amount of data of the object's surface can never be perfect, as the resolution is arbitrary, thus leading to many injective reconstruction methods.

One of the earliest, but still well-known algorithms, presented by
 Lorensen et al. \cite{Lorensen:1987:MCH:37402.37422} reconstructs the surface of an object as a
  triangle mesh, given three-dimensional medical data. This data mostly is given as
   a three-dimensional discrete scalar field, also known as voxels. 
During the reconstruction process, the algorithm iterates through all voxels. For each of these positions, up to eight occupied positions of the voxel-grid can be considered. Furthermore, a reconstruction configuration is defined for every configuration of present/non-present positions in the neighborhood of the voxel. Thus, leading to $2^8 = (256)$ possible configurations, where 15 of them are unique, and the rest are generated by the rotation of the base configurations.  Every configuration reconstructs a number of polygons given one voxel configuration. Finally, all polygons created this way are merged into one mesh.

Similar to the marching cubes algorithm, the work presented by 
Bernardini et al. \cite{817351}, the ball-pivoting algorithm, 
interpolates the surface of an object given the input data.
In contrast, the algorithm processes a collection of unstructured 
three-dimensional points in space with their orientation on the
(ground truth) surface and returns a triangle mesh. Starting from
 a seed triangle, for each edge of not already processed triangles
 , a sphere with set radius $r$ revolves around it. If another point
  of the input point cloud intersects with the sphere, a triangle is 
  created out of the endpoints of the considered edge and the newly 
  found point, thus creating a triangle mesh if every edge of each 
  triangle has been considered. 
  As a baseline, results by the process presented in this work
  will be compared to those of the results by the ball-pivoting algorithm.
\begin{enumerate}
  \item poisson surface reconstruction
  \item Fabians \cite{Groh2017}
  \item instant field aligned meshes
  \item Dennis w\cite{bukenberger2018hierarchical}
\end{enumerate}


\section{Machine Learning review}
\todo{different name for section}
\label{ml_review}

\section{Machine Learning based approaches}
\label{ml_approaches}
\todo{different name for section}
\begin{enumerate}
  \item machine learning good way for inference, probably neural network too, given huge amount of data and finding similarities in data
  \item many approaches for surface reconstruction in classic ml
  \item used for self driving cars. Fast solutions
  \item NN recently started to get nice results
  \item many try to transform given input data to voxel based representation
  \item not many directly from point cloud to meshes
  \item range scanner to meshes
  \item end result not meshes?
\end{enumerate}
Papers to cite:
\begin{enumerate}
  \item Convolutional neural network
  \item Semi-Supervised Classification with Graph Convolutional networks
  \item dense 3d object reconstruction from single depth view
  \item PointNet++
  \item deep marching cubes
  \item pixel2mesh
  \item learning a hierarchical latent-variable model of 3d shapes
  \item FlexConv
  \item unsupervised learning of 3d structure
  \item image2 mesh
  \item Surface reconstruction from unorganized Points
\end{enumerate}
\todo{Find non NN ML papers, from the other prof of ML lecture?}
\todo{Saliency for feature detection}


x`\chapter{Related Work}
\label{sec:relatedwork}
Often in computer graphics, it is necessary to process three-dimensional objects. There are many ways to acquire data of the surface of such geometry and more so many techniques to transform and augment that data to distinct representations. This often non-trivial task is crucial to further process the object in question in later stages of their respective pipeline. Over the years many representations of acquisition data formats and transformation methods, as well as target data formats, have accumulated.\\
In this chapter, various of these techniques and data formats are examined, of which some of them are used in this work as a vehicle for a novel data transformation routine.\\
Initially, classical approaches are examined in section \ref{classic_approaches} which do not rely on artificial intelligence or machine learning methods.
With the recent advances in deep learning, many new approaches have been developed and thus considered in this work.
Subsequently, machine learning based methods are reassessed in section \ref{ml_approaches} which rely on distinctive statistical features in their initial data format or the dataset itself, thus allowing for the transformation.
\todo{citations for this stuff?}
\section{Classical approaches}
\label{classic_approaches}
In this chapter, classical approaches are examined which do not rely on machine
 learning to transform data of objects from one format to another. 

In general, a three-dimensional object may be specified by its representation of the
 surface. This representation varies from unstructured point cloud data in three-dimensional space
  to a graph-based representation like triangle-/ or quad-meshes and even poly meshes
   with more than four neighbors per node. Information on the object's surface may also be described
    by three-dimensional discrete scalar field (~Voxels~).  Often the information on an object's
     surface is only partially defined, where some information may be missing. 
     This can be due to acquisition methods, where parts of the object are partially 
     concealed or recorded from insufficient many sides, or that information is not available.

Starting from these representations of an object's surface, there are many methods 
of transformation to reach another state of representation. 
Transforming the data not only makes further processing easier
 but rather during their process, more data on the surface is 
 computed based on the given input data.

 This task is non-trivial since, from one representation of an object's surface, many reconstructions in the form of another representation are possible. The amount of data of the object's surface can never be perfect, as the resolution is arbitrary, thus leading to many injective reconstruction methods.
 In this section, some of the more important works on this subject of data transformation are examined, with 
 an overview of the different input and output formats available. 

 \subsection{Transformation methods}
 \todo{maybe not needed. only classic approaches as section?}

One of the earliest, but still well-known algorithms, presented by
 Lorensen et al. \cite{Lorensen:1987:MCH:37402.37422}, the marching cubes algorithm, reconstructs the surface information of an object as a
  triangle mesh, given three-dimensional medical data. This data is provided in the form of
   voxels. 
During the reconstruction process, the algorithm iterates through all voxels. For each of these positions, up to eight occupied positions of the voxel-grid can be considered. Furthermore, a reconstruction configuration is defined for every configuration of present/non-present positions in the neighborhood of the voxel. Thus, leading to $2^8 = (256)$ possible configurations, where 15 of them are unique, and the rest are generated by the rotation of the base configurations.  Every configuration reconstructs a number of polygons given one voxel configuration. Finally, all polygons created this way are merged into one mesh.
 
A more recent work by Groh et al.\cite{Groh2017} presents a way to digitalize real-world objects by projecting shifted sinusoid patterns by multiple projectors onto the object. This is captured by multiple cameras. By combining the shifted patterns and highly accurate created features, a precise bundle adjustment is performed. Thus, all projector-camera pairs for each surface point can be estimated and used to optimize the final depth information.
\todo{Actually not sure if this fits here}

By all means, digitalizing real-world objects not only revolves on structured light but also may make use of range scanners, which measure the depth from the recording device to a point on the surface of the object. Additionally, this depth information may be combined from recording from multiple views. Assembling this information yields unstructured three-dimensional point cloud data. 

Similar to the marching cubes algorithm, the work presented by 
  Bernardini et al. \cite{817351}, the ball-pivoting algorithm, 
  interpolates the surface of an object given the input data.
  In contrast, the algorithm processes such a collection of unstructured 
  three-dimensional points in space with their orientation on the
  ground truth surface and returns a triangle mesh. Starting from
   a seed triangle, for each edge of not already processed triangles
   , a sphere with set radius $r$ revolves around it. If another point
    of the input point cloud intersects with the sphere, a triangle is 
    created out of the endpoints of the considered edge and the newly 
    found point, thus creating a triangle mesh if every edge of each 
    triangle has been considered. 
  

    Furthermore, the Poisson surface reconstruction, 
    presented by Kazhdan et al. \cite{Kazhdan:2006:PSR:1281957.1281965}, also processes a
     collection of unstructured points in three-dimensional
      space with surface orientation. Solving the Poisson problem 
      \cite{Genovese2006} assuming the isosurface of the surface can be 
      approximated by the normal field, thus yielding the
       indicator function. This function results in non-zero
        values close to the surface of the object and thereby 
        can be directly extracted. The problem is then reduced to a 
        form of voxel-grid, thus can be transformed by the marching
         cubes algorithm.

 Similarly, surface information can be reconstructed by calculating 
 the PC's medial axis transform ~(~MAT~) from unstructured point cloud
  data, followed by an inverse transformation.  The medial axis of a PC
   is described as a set of points where each point is the closest neighbor 
   to at least one other point in the original PC. Not only can the PC be noisy 
   or scanty, as described by Amenta et al. \cite{Groh2017}, but the reconstructed
    surface mesh will always be manifold.
    
  
  
\begin{enumerate}
  \item surface reconstruction from unorganized points
  \item instant field aligned meshes
  \item Dennis \cite{bukenberger2018hierarchical}
\end{enumerate}


\section{Machine Learning based approaches}
\label{ml_approaches}
\todo{different name for section}
\begin{enumerate}
  \item machine learning good way for inference, probably neural network too, given huge amount of data and finding similarities in data
  \item many approaches for surface reconstruction in classic ml
  \item used for self driving cars. Fast solutions
  \item NN recently started to get nice results
  \item many try to transform given input data to voxel based representation
  \item not many directly from point cloud to meshes
  \item range scanner to meshes
  \item end result not meshes?
\end{enumerate}
Papers to cite:
\begin{enumerate}
  \item Convolutional neural network
  \item Semi-Supervised Classification with Graph Convolutional networks
  \item dense 3d object reconstruction from single depth view
  \item PointNet++
  \item deep marching cubes
  \item pixel2mesh
  \item learning a hierarchical latent-variable model of 3d shapes
  \item FlexConv
  \item unsupervised learning of 3d structure
  \item image2 mesh
  \item Surface reconstruction from unorganized Points
\end{enumerate}
\todo{Find non NN ML papers, from the other prof of ML lecture?}
\todo{Saliency for feature detection}


\chapter{Material}
\label{sec:material}

\section{Data}
Explain and show 

\chapter{Methods}
\label{sec:methods}

\chapter{Results}
\label{chap:results}

    By the nature of the defined configurations $\mathcal{C}_i$, the neural network
    \emph{points2mesh} has several reconstructions for each input point cloud $\mathcal{PC}$.
    Since it is not easy to estimate the quality of reconstructed meshes only based on numerical
    values, it is vital to assess their visual quality, tough highly subjective.
    Therefore, in this chapter, a wide range of reconstructed objects of various categories $\mathcal{C}_i$
    are presented and displayed next to their input point cloud $\mathcal{PC}$ and their ground 
    truth mesh. Some interesting features of reconstructions are pointed out, and presented with reconstructions of comparable
    reconstruction methods, as well as the alternative approach to train the network as outline in \ref{trainings}.
    
\section{Compared methods}

Restoring polygonal meshes directly from point cloud data is a new domain of geometry reconstruction, 
which yet has to be explored deeply. Apart from \emph{points2mesh}, no other methods relying on neural 
networks and deep learning techniques have yet been proposed for this kind of mode of data transformation.
For this reason, \emph{points2mesh} has to compete against techniques from traditional approaches like
\emph{ball-pivoting algorithm} \cite{817351}(~BPA~) and \emph{instant field meshes} \cite{Jakob2015Instant} (~IFM~). Furthermore, by adding
an intermediate step of voxelization of $\mathcal{PC}$, this work also can be compared against 
\emph{deep marching cubes} \cite{Liao2018CVPR} (~DMC~), which transforms voxelized data into meshes based on deep learning approaches and the classical
marching cubes algorithm.

\section{Result visualization}
% 1024Figure
  \ref{fig:1} shows the reconstruction from a point cloud sample size of 1024 from the \emph{category}
  \emph{airplane}. \ref{fig:1a} shows the ground truth mesh, and \ref{fig:1d} its sampled point cloud. Figures
  \ref{fig:1b}, \ref{fig:1c}, \ref{fig:1e}, \ref{fig:1f} show the reconstruction from $\mathcal{N}_{recon}$ with 
  configurations $\mathcal{C}_1$ through $\mathcal{C}_4$.

\begin{figure}[htbp]
  \centering
  \subfloat[Ground truth mesh.\label{fig:1a}]{\includegraphics[width=0.32\textwidth]{gt_airplane_1.png}}
  \subfloat[Configuration $\mathcal{C}_1$.\label{fig:1b}]{\includegraphics[width=0.32\textwidth]{c1_1024_airplane_1.png}}
  \subfloat[Configuration $\mathcal{C}_2$.\label{fig:1c}]{\includegraphics[width=0.32\textwidth]{c2_1024_airplane_1.png}}\\
  \subfloat[Sampled point cloud from ground truth mesh.\label{fig:1d}]{\includegraphics[width=0.32\textwidth]{airplane_1_pc.png}}
  \subfloat[Configuration $\mathcal{C}_3$.\label{fig:1e}]{\includegraphics[width=0.32\textwidth]{c3_1024_airplane_1.png}}
  \subfloat[Configuration $\mathcal{C}_4$.\label{fig:1f}]{\includegraphics[width=0.32\textwidth]{c4_1024_airplane_1.png}}\\
  \subfloat[Reconstrucion with \emph{BPA}.\label{fig:1g}]{\includegraphics[width=0.32\textwidth]{airplane_bpa_1024.png}}
  \subfloat[Reconstruction with \emph{IFM}.\label{fig:1h}]{\includegraphics[width=0.32\textwidth]{ifm_1024_airplane_1.png}}
  \subfloat[Reconstruction with \emph{DMC}.\label{fig:1i}]{\includegraphics[width=0.32\textwidth]{dmc_airplane.png}}
  \caption{Reconstruction of an airplane from point cloud with 1024 samples.} \label{fig:1}
\end{figure}

\begin{figure}[htbp]
  \centering
  \subfloat[Ground truth mesh.\label{fig:2a}]{\includegraphics[width=0.32\textwidth]{gt_chair_1.png}}
  \subfloat[Configuration $\mathcal{C}_1$.\label{fig:2b}]{\includegraphics[width=0.32\textwidth]{c1_1024_chair_1.png}}
  \subfloat[Configuration $\mathcal{C}_2$.\label{fig:2c}]{\includegraphics[width=0.32\textwidth]{c2_1024_chair_1.png}}\\
  \subfloat[Configuration $\mathcal{C}_3$.\label{fig:2d}]{\includegraphics[width=0.32\textwidth]{c3_1024_chair_1.png}}
  \subfloat[Configuration $\mathcal{C}_4$.\label{fig:2epro}]{\includegraphics[width=0.32\textwidth]{c4_1024_chair_1.png}}
  \caption{Reconstruction of a chair from point cloud with 1024 samples.} \label{fig:2}
\end{figure}

\begin{figure}[htbp]
  \centering
  \subfloat[Ground truth mesh.\label{fig:3a}]{\includegraphics[width=0.32\textwidth]{gt_guitar_1.png}}
  \subfloat[Configuration $\mathcal{C}_1$.\label{fig:3b}]{\includegraphics[width=0.32\textwidth]{c1_1024_guitar_1.png}}
  \subfloat[Configuration $\mathcal{C}_2$.\label{fig:3c}]{\includegraphics[width=0.32\textwidth]{c2_1024_guitar_1.png}}\\
  \subfloat[Sampled point cloud from ground truth mesh.\label{fig:3d}]{\includegraphics[width=0.32\textwidth]{guitar_pc00.png}}
  \subfloat[Configuration $\mathcal{C}_3$.\label{fig:3e}]{\includegraphics[width=0.32\textwidth]{c3_1024_guitar_1.png}}
  \subfloat[Configuration $\mathcal{C}_4$.\label{fig:3f}]{\includegraphics[width=0.32\textwidth]{c4_1024_guitar_1.png}}\\
  \subfloat[\emph{IFM} guitar reconstruction.\label{fig:3g}]{\includegraphics[width=0.32\textwidth]{ifm_1024_guitar_1.png}}
  \subfloat[\emph{DMC} guitar reconstruction.\label{fig:3h}]{\includegraphics[width=0.32\textwidth]{dmc_guitar.png}}
  \caption{Reconstruction of a guitar from point cloud with 1024 samples.} \label{fig:3}
\end{figure}
\begin{figure}[htbp]
  \centering
  \subfloat[Ground truth mesh.\label{fig:4a}]{\includegraphics[width=0.32\textwidth]{gt_bowl_1.png}}
  \subfloat[Configuration $\mathcal{C}_1$.\label{fig:4b}]{\includegraphics[width=0.32\textwidth]{c1_1024_bowl_1.png}}
  \subfloat[Configuration $\mathcal{C}_2$.\label{fig:4c}]{\includegraphics[width=0.32\textwidth]{c2_1024_bowl_1.png}}\\
  \subfloat[Reconstruction with \emph{DMC}\label{fig:4d}]{\includegraphics[width=0.32\textwidth]{dmc_bowl.png}}
  \subfloat[Configuration $\mathcal{C}_3$.\label{fig:4e}]{\includegraphics[width=0.32\textwidth]{c3_1024_bowl_1.png}}
  \subfloat[Configuration $\mathcal{C}_4$.\label{fig:4f}]{\includegraphics[width=0.32\textwidth]{c4_1024_bowl_1.png}}
  \caption{Reconstruction of a bowl from point cloud with 1024 samples.} \label{fig:4}
\end{figure}

  The last row depicts reconstructions with \emph{BPA}, \emph{IFM}, and \emph{DMC}. While the reconstruction with configuration 
  $\mathcal{C}_1$ represents the most interesting features of the airplane, especially the small winglets, it still has some errors. 
  One of the errors is at the rear of the plane with an edge dangling in the air, spanning over it. The same error happens with configuration
  $\mathcal{C}_3$ and $\mathcal{C}_4$ too. Configuration $\mathcal{C}_2$ seems to represent the airplane's features the best, though the least
  smooth of all reconstructions. As seen in figures \ref{fig:1f} and \ref{fig:2c}, configuration $\mathcal{C}_2$ is generally depicting the ground
    truth mesh in a good way, though not smooth.
  Considering other reconstruction methods, different problems, but also positives can be seen. \emph{BPA} reproduces smooth results as seen 
  in figure \ref{fig:1g}, while \emph{IFM} does not perform as well on low samples (~256 and 1024~) as seen in figure \ref{fig:1h}. Generally 
  \emph{DMC} does not perform good either as seen in figures \ref{fig:1i} and \ref{fig:3h}.

  Similarly, figure \ref{fig:2} through \ref{fig:5} show reconstructions with the same configurations $\mathcal{C}_i$ from $\mathcal{N}_{recon}$,
  but with an object from the categories \emph{chair}, \emph{guitar}, \emph{bowl}. These categories were included
  in the training set of $\mathcal{N}_{recon}$, though every reconstructed object sample originate from a test set, and thus have not be seen
  during the training stage.
  Similar behavior of the four network configurations, as described above, can be seen again, including the spanning of edges over empty space,
  like in figure \ref{fig:2d}, between the legs of the chair.


%% 256
%\begin{figure}[htbp]
%  \centering
%  \centering
%  \subfloat[Ground truth mesh.\label{fig:5a}]{\includegraphics[width=0.5\textwidth]{gt_chair_256_1.png}}
%  \subfloat[Ground truth mesh.\label{fig:5b}]{\includegraphics[width=0.5\textwidth]{gt_chair_7500.png}}\\
%  \subfloat[Sampled point cloud from ground truth mesh.\label{fig:5b}]{\includegraphics[width=0.5\textwidth]{chair_256_pc06.png}}
%  \subfloat[Sampled point cloud from ground truth mesh.\label{fig:5c}]{\includegraphics[width=0.5\textwidth]{chair_256_pc00.png}}\\
%  \subfloat[Configuration $\mathcal{C}_1$.\label{fig:5d}]{\includegraphics[width=0.5\textwidth]{c1_256_chair_1.png}}
%  \subfloat[Configuration $\mathcal{C}_2$.\label{fig:5e}]{\includegraphics[width=0.5\textwidth]{c1_256_chair_2.png}}\\
%  \caption{Reconstruction of chairs from point cloud with 256 samples.} \label{fig:5}
%\end{figure}
%\begin{figure}[htbp]
%  \centering
%  \subfloat[Ground truth mesh.\label{fig:51a}]{\includegraphics[width=0.5\textwidth]{gt_car_1.png}}
%  \subfloat[Sampled point cloud from ground truth mesh.\label{fig:51b}]{\includegraphics[width=0.5\textwidth]{car_pc05.png}}\\
%  \subfloat[Configuration $\mathcal{C}_1$.\label{fig:51c}]{\includegraphics[width=0.5\textwidth]{c1_256_car_1.png}}
%  \subfloat[Reconstruction of car with \emph{DMC},.\label{fig:51d}]{\includegraphics[width=0.5\textwidth]{dmc_car.png}}
%  \caption{Reconstruction of a car from point cloud with 256 samples.} \label{fig:51}
%\end{figure}
%\begin{figure}[htbp]
%  \centering
%  \subfloat[Ground truth mesh.\label{fig:52a}]{\includegraphics[width=0.33\textwidth]{gt_256_toilet.png}}
%  \subfloat[Configuration $\mathcal{C}_1$.\label{fig:52b}]{\includegraphics[width=0.33\textwidth]{c1_256_toilet.png}}
%  \subfloat[Reconstruction with \emph{BPA}\label{fig:52c}]{\includegraphics[width=0.33\textwidth]{bpa_256_toilet.png}}
%  \caption{Reconstruction of a toilet from point cloud with 256 samples.} \label{fig:52}
%\end{figure}
%  Figures \ref{fig:5} and \ref{fig:51} show reconstruction with a sample size of 256. Depicted are reconstructions with configurations
%  $\mathcal{C}_1$ and $\mathcal{C}_2$ 
%  and their corresponding ground truth mesh and sampled point cloud.
%  While still able to reconstruct the general shape of the object, the spanning problem is still visible.
%  Additionally, \emph{DMC} is not able to handle bigger empty spaces in the sparse cloud well enough to create a proper reconstruction,
%  as seen in figure \ref{fig:51d}.
%
%
%  \begin{figure}[htbp]
%    \centering
%    \subfloat[Ground truth mesh.\label{fig:6a}]{\includegraphics[width=0.4\textwidth]{gt_chair_7500.png}}
%    \subfloat[Sampled point cloud from ground truth mesh.\label{fig:6b}]{\includegraphics[width=0.4\textwidth]{chair_pc00.png}}\\
%    \subfloat[Configuration $\mathcal{C}_1$.\label{fig:6c}]{\includegraphics[width=0.4\textwidth]{c1_7500_chair_2.png}}
%    \subfloat[Configuration $\mathcal{C}_4$.\label{fig:6d}]{\includegraphics[width=0.4\textwidth]{c4_7500_chair_2.png}}
%    \caption{Reconstruction of a chair from point cloud with 7500 samples.} \label{fig:6}
%  \end{figure}
%  \begin{figure}[htbp]
%    \centering
%    \subfloat[Ground truth mesh.\label{fig:61a}]{\includegraphics[width=0.4\textwidth]{gt_airplane_7500.png}}
%    \subfloat[Samples pointc loud from ground truth mesh.\label{fig:6d}]{\includegraphics[width=0.4\textwidth]{plane2_pc04.png}}\\
%    \subfloat[Configuration $\mathcal{C}_1$.\label{fig:6e}]{\includegraphics[width=0.4\textwidth]{c4_7500_airplane_2.png}}
%    \subfloat[Configuration $\mathcal{C}_4$.\label{fig:6f}]{\includegraphics[width=0.4\textwidth]{c1_7500_airplane_2.png}}
%    \caption{Reconstruction of an airplane from point cloud with 7500 samples.} \label{fig:61}
%  \end{figure}
%  
%  Figures \ref{fig:6} and \ref{fig:61} show reconstruction of a chair and an airplane from 7500 samples. Depicted are configurations 
%  $\mathcal{C}_1$ and $\mathcal{C}_4$.
%
%
%% 7500
%
%%256-1024-7500
%\begin{figure}[htbp]
%  \centering
%  \subfloat[Ground truth mesh.\label{fig:7a}]{\includegraphics[width=0.5\textwidth]{gt_airplane_7500.png}}
%  \subfloat[$\mathcal{C}_1$ with sample size 256.\label{fig:7b}]{\includegraphics[width=0.5\textwidth]{c1_256_airplane_2.png}}\\
%  \subfloat[$\mathcal{C}_1$ with sample size 1024.\label{fig:7c}]{\includegraphics[width=0.5\textwidth]{c1_1024_airplane_2.png}}
%  \subfloat[$\mathcal{C}_1$ with sample size 7500.\label{fig:7d}]{\includegraphics[width=0.5\textwidth]{c4_7500_airplane_2.png}}\\
%  \subfloat[\emph{BPA} with sample size 7500.\label{fig:7e}]{\includegraphics[width=0.5\textwidth]{airplane2_ifm_7500.png}}
%  \subfloat[\emph{IFM} with sample size 7500.\label{fig:7f}]{\includegraphics[width=0.5\textwidth]{airplane2_bpa_7500.png}}
%  \caption{Reconstruction of an airplane with different sample sizes, but the same configuration $\mathcal{C}_1$.} \label{fig:7}
%\end{figure}
%
%  For easier comparisons, figure \ref{fig:7} shows the reconstruction of the same airplane with configuration $\mathcal{C}_1$ but with
%  different sample sizes, as well as reconstructions from \emph{BPA} and \emph{IFM} with the same sample size, showing their ostensible
%  superiority with more samples.
%
%
%  \begin{figure}[htbp]
%    \centering
%    \subfloat[Noisy airplane 1.\label{fig:81a}]{\includegraphics[width=0.5\textwidth]{c1_1024_noise_airplane_1.png}}
%    \subfloat[Noisy airplane 2.\label{fig:81b}]{\includegraphics[width=0.5\textwidth]{c1_1024_noise_airplane_2.png}}\\
%    \subfloat[Noisy airplane 3.\label{fig:81c}]{\includegraphics[width=0.5\textwidth]{c1_1024_noise_airplane_3.png}}
%    \subfloat[Noisy airplane 4.\label{fig:81d}]{\includegraphics[width=0.5\textwidth]{c1_1024_noise_airplane_4.png}}
%    \caption{Reconstruction of airplanes with $\mathcal{C}_1$ from noisy point cloud with noise factor of $0.01$. } \label{fig:81}
%  \end{figure}
%  Next, in figure \ref{fig:81}, reconstruction of airplanes and with configuration $\mathcal{C}_1$ can be seen. Each point in the point 
%  cloud has random
%  jitter noise between $-0.01$ and $0.01$ added on top of it.
%
%
%% generalization
%\begin{figure}[htbp]
%  \centering
%  \subfloat[Ground truth mesh.\label{fig:8a}]{\includegraphics[width=0.32\textwidth]{gt_person_1.png}}
%  \subfloat[$\mathcal{C}_1$.\label{fig:8b}]{\includegraphics[width=0.32\textwidth]{c1_1024_person_1.png}}
%  \subfloat[$\mathcal{C}_2$.\label{fig:8c}]{\includegraphics[width=0.32\textwidth]{c2_1024_person_1.png}}\\
%  \subfloat[Sampled point cloud from ground truth mesh.\label{fig:8d}]{\includegraphics[width=0.32\textwidth]{person_pc00.png}}
%  \subfloat[$\mathcal{C}_3$.\label{fig:8e}]{\includegraphics[width=0.32\textwidth]{c3_1024_person_1.png}}
%  \subfloat[$\mathcal{C}_4$.\label{fig:8f}]{\includegraphics[width=0.32\textwidth]{c4_1024_person_1.png}}
%  \caption{Reconstruction from point clouds of untrained category \emph{person}.} \label{fig:8}
%\end{figure}
%\begin{figure}[htbp]
%  \centering
%  \subfloat[$\mathcal{C}_1$.\label{fig:9a}]{\includegraphics[width=0.5\textwidth]{c1_1024_bathtub_1.png}}
%  \subfloat[$\mathcal{C}_2$.\label{fig:9b}]{\includegraphics[width=0.5\textwidth]{c2_1024_bathtub_1.png}}\\
%  \subfloat[$\mathcal{C}_3$.\label{fig:9c}]{\includegraphics[width=0.5\textwidth]{c3_1024_bathtub_1.png}}
%  \subfloat[$\mathcal{C}_4$.\label{fig:9d}]{\includegraphics[width=0.5\textwidth]{c4_1024_bathtub_1.png}}
%  \caption{Reconstruction from point clouds of untrained category \emph{bathutb}.} \label{fig:9}
%\end{figure}
%  Additionally, figure \ref{fig:8} and figure \ref{fig:9} show the reconstruction of a person and a bathtub, based on point cloud samples 
%  which $\mathcal{N}_{recon}$ has not been trained with.
%\begin{figure}[htbp]
% \centering
% \subfloat[Reconstruction with 256 samples\label{fig:12a}]{\includegraphics[width=0.33\textwidth]{c1_opt_256_airplane2.png}}
%  \subfloat[Reconstruction with 1024 samples\label{fig:12b}]{\includegraphics[width=0.33\textwidth]{c1_opt_1024_airplane2.png}}
%  \subfloat[Reconstruction with 7500 samples\label{fig:12c}]{\includegraphics[width=0.33\textwidth]{c1_opt_7500_airplane2.png}}
%  \caption{Reconstruction from point clouds of airplanes of $\mathcal{N}_{opt}$. Input point cloud is equal to ground truth point cloud} \label{fig:10}
%\end{figure}
%\begin{figure}[htbp]
%  \centering
%  \subfloat[\label{fig:11a}]{\includegraphics[width=0.5\textwidth]{c1_opt_1024_guitar.png}}
%  \caption{Reconstruction from point clouds of guitar of $\mathcal{N}_{opt}$. Input point cloud is equal to ground truth point cloud with 1024 samples} \label{fig:11}
%\end{figure}
%\begin{figure}[htbp]
%  \centering
%  \subfloat[\label{fig:10a}]{\includegraphics[width=0.3\textwidth]{c1_opt_1024_chair1.png}}
%  \subfloat[\label{fig:10b}]{\includegraphics[width=0.3\textwidth]{c1_opt_1024_chair2.png}}
%  \subfloat[\label{fig:10c}]{\includegraphics[width=0.3\textwidth]{c1_opt_1024_chair3.png}}
%  \caption{Reconstruction from point clouds of chairs of $\mathcal{N}_{opt}$. Input point cloud is equal to ground truth point cloud with 1024 samples} \label{fig:10}
%\end{figure}
%
  Finally, figure \ref{fig:10} shows the reconstruction of chairs, based on $\mathcal{N}_{opt}$. Only the reconstruction with the configuration $\mathcal{C}_1$ is shown
  since it produces the most interesting results\footnote{See Appendix \ref{optrecon} for some reconstruction with other configurations}. Configurations $\mathcal{C}_2$ through $\mathcal{C}_4$ produce similar results. Though no ground truth data was available during training time, reconstruction
  of the airplane is still possible. It is visible, that only 256 samples are not enough to create a good reconstruction. Starting with 1024 samples, the airplane reconstruction
  is closer to the same reconstruction as with $\mathcal{N}_{recon}$ as seen in \ref{fig:81}. Similarly, as seen in figure \ref{fig:11}, the reconstructed guitar can be 
  seen, though not as clean as the guitar from \ref{fig:3}. Reconstrucions of different chairs as seen in figure \ref{fig:10} with $\mathcal{N}_{opt}$ also show the problem with dangling edges in empty space
  to be a more common occurence. 


  Reconstructing watertight meshes allows to easily 3D print them. Figure \ref{fig:13} shows the reconstruction with $\mathcal{N}_{recon}$ with configuration $\mathcal{C}_1$. 
  The prediction then is fed into a 3D toolpath generator (~slicer~)\cite{Ranellucci2011} and then printed with a \emph{PRUSA I3 MK3S} at $0.015mm$ layer hight, taking approximately 8 hours printing time.
  \begin{figure}[htbp]
    \centering
    \subfloat[3D printed guitar reconstructed with $\mathcal{N}_{recon}$\label{fig:13a}]{\includegraphics[width=0.3\textwidth]{guitar_print.png}}
    \subfloat[3D printed car reconstructed with $\mathcal{N}_{recon}$\label{fig:13b}]{\includegraphics[width=0.3\textwidth]{car_print.png}}
    \subfloat[3D printed toilet reconstructed with $\mathcal{N}_{recon}$\label{fig:13c}]{\includegraphics[width=0.3\textwidth]{toilet_print.png}}
    \caption{3D printed reconstruction of guitar, car and toilet of $\mathcal{N}_{recon}$ and $\mathcal{C}_1$ with 1024 samples.} \label{fig:13}
  \end{figure}


\chapter{Evaluation}
\label{chap:evaluation}

    An indefinite amount of different meshes can represent the same object.
    Consequently, even more meshes exist, which approximate that same object. 
    Therefore, quantifying the quality of reconstructed meshes is a critical aspect for 
    this application, likewise for similar ones. Only by being able to determine applications
    with higher qualitative results facilitate improvement in the field of mesh 
    reconstruction. 
    \emph{points2mesh} and its described configurations $\mathcal{C}_i$ aim to improve on previous work,
    reconstructing meshes from point cloud data. Its advantages and drawbacks, in contrast to similar efforts, 
    or even reconstruction methods working on other data types, rather than point cloud data, have to be clearly
    defined and evaluated.
    Even though the goal of this work is to reconstruct visually pleasing reconstructions for mesh approximations, 
    considering raw numeric evaluations is nevertheless a crucial component.

    In section \ref{subsec:exp}, the experimental setup is described, which tries to assess
    the quality of reconstructed meshes and compare the defined metrics to other reconstruction methods. 
    Subsequently, the results of the experiment are presented in section \ref{subsec:results}, where
    each one is accompanied by tabulated results.


\section{Experimental setup}
\label{subsec:exp}

    Finding a proper way to quantify the approximation of a mesh in comparison to its reconstruction is no straightforward task.
    Many different methods exist, each with its advantages and drawbacks. Depending on the specific application like the size of
    the mesh, the number of vertices and edges, or features to compare, optimal methods of comparisons may differ, if any even
    exist. For the proposed network \emph{points2mesh}, several metrics are chosen, which may lead to conclusions about the quality 
    of the reconstructed mesh. Not only are they applicable to \emph{points2mesh}, but to other reconstruction methods as well. 
    Thus, these metrics can be compared to each other, and an argument may be formed about their reconstruction capacity. This 
    section specifies aspects to evaluate of \emph{points2mesh}'s configurations $\mathcal{C}_i$. Furthermore, utilized metrics 
    for the comparison are detailed in the subsequent subsection. Finally, competing methods for reconstruction are specified,
    to define a frame of reference against which \emph{points2mesh} should be able to hold up or even exceed.
    Finally, a collection of inference examples of the second approach for network training as detailed in \ref{trainings} are listed as well.


\subsection{Evaluation aspects}
    There are several aspects to evaluate \emph{points2mesh}. Not only are network configurations $\mathcal{C}_i$ a crucial consideration,
    but so is computation time, as well as generalization capabilities and starting meshes $\mathcal{M}_i$ (~Ellipsoid, Torus~). 
    Finally,
    attaining visually pleasing meshes is one such goal too, and thus has to be evaluated as well.

    Therefore, the following aspects are chosen to evaluate.
    \begin{itemize}
        \item \textbf{Configurations} $\mathcal{C}_i$: Each configuration, as described in \ref{configurations}, has its purpose 
                for specific reconstruction capabilities. Evaluating every one of them may lead to conlcusions about objective rankings, and thus
                to a choice of best network configuration for reconstructions with the highest quality.
        \item \textbf{Input sample size}: Based on the input sample size of the point cloud, harder, or easier tasks for deforming the ellipsoid are assumed. 
                Evaluating the quality of the reconstructions may lead to conclusions about the amount of necessary samples for proper inference.
                The input sample sizes are $\in[256,1024,7500]$. 
        \item \textbf{Visual quality}: While being subjective, taking a closer look how good the reconstruction seems to represent the ground truth mesh
        is still relevant.
        \item \textbf{Generalization capabilities}: Evaluating how well configrations $\mathcal{C}_i$ are capable of generalize the problem of reconstruction is a crucial aspect
        for any neural network. 
        \item \textbf{Changing the basemesh $\mathcal{M}_i$}: \emph{points2mesh} is only able to reconstruct meshes of the same \emph{genus} as the input mesh $\mathcal{M}_i$.
        Thus, testing configurations $\mathcal{C}_i$ with different $\mathcal{M}_i$ of higher \emph{genus} like a toroidal shape may lead to different results and 
        reconstructions with higher \emph{genus}.
        \item \textbf{Noise robustness}: Acquiring unstructured points in three-dimensional space from real-world object often entails noisy data.
         The perfect acquisition is already nearly impossible due to imperfection in the hardware. Thus, training $\mathcal{N}_{recon}$ with noisy data 
         is an essential part of the evaluation.
        \item \textbf{Computation time}: When reconstruction methods are used for time critical application, their computation time is important to evaluate.
    \end{itemize}

    While evaluating configurations, the visual quality and computation time can be compared to other methods, generalization and different used basemeshes
    are isolated during evaluation due to their nature and thus are only compared to each other.

\subsection{Evaluation metrics}

    Choosing proper evaluation metrics is ambiguous, but several are often chosen in recent work of mesh reconstruction with favorable information content. 
    While some of them may help to lead to objective conclusions about quality, others are subjective in nature. These metrics 
    are described in the following. 

    \subsubsection*{Meshdistance $d_m$}    
    In 3D vision, a commonly used metric to compare the quality of meshes is done with completeness of the
    predicted mesh in comparison to the ground truth mesh. Based on uniformly sampled points on both meshes, completeness describes the distance of points of
    ground truth to the predicted mesh. 
    Stutz et al. \cite{Stutz2017, Stutz2018CVPR} introduce a parallel $C++$ implementation of \emph{meshdistance} which
    calculates completeness, given a ground truth mesh and a prediction. Distances are calculated by a
    point-to-triangle method by Christopher Batty\footnote{https://github.com/christopherbatty/SDFGen last visited: \today} (~See chapter \ref{back:ptt})
    %\subsubsection*{Chamfer distance $d_C$}
    %Similarly to the \emph{meshdistance} method, to calculate the \emph{chamfer distance}, both meshes have to be uniformly sampled. 
    %In contrast, only the distance to each nearest neighbor from ground truth to reconstruction, as well as from reconstruction 
    %to ground truth, is considered for the metric.
    %While the \emph{chamfer distance} is a more simplified version of the meshdistance, it is still able to provide
    %valuable information about the quality of reconstructions.
    %\subsubsection*{MAYBE Earthmovers distance}
    %\subsubsection*{Hausdorff distance $d_H$}
    %The \emph{hausdorff distance} describes a more sophisticated, asymmetric metric for measuring distances between two similar meshes. 
    %Anew, the ground truth mesh is sampled with the same amount of samples as the ground truth mesh possesses. For each sample, its 
    %closest point on the predicted mesh is calculated, resulting in minimum, maximum values, as well as mean values for the \emph{hausdorff distance} metric.
    %\begin{align}
    %    d_{H}(\mathcal{M}_{gt},\mathcal{M}_i) = \underset{y\in\mathcal{M}_i}{sup} \underset{y\in\mathcal{M}_{gt}}{inf} d(x,y)
    %\end{align}
    %With $d(x,y)$ describing a function, calculating the distance from $x$ to $y$.
    \subsubsection*{Visual evidence}    
    Though subective, taking a closer look at, how well a reconstructed mesh seems to represent the ground truth mesh is important too. Relying only on 
    numeric values may lead to reconstructions with low values, but not visually pleasing results. This is due to the complicated nature of the problem itself.
    For a given point cloud, an indefinite amount of different reconstructions are possible. Therefore, inspecting the meshes is necessary.


\section{Experiment results}
\label{subsec:results}
    The numerical evaluation of \emph{points2mesh} and its configuration $\mathcal{C}_i$ is denoted in a tabulated form.
    Since \emph{points2mesh} is trained on the ModelNet dataset, data samples from this dataset are selected for inference. About 20 percent 
    of the provided point clouds are left out of training and thus used as the test-set. Therefore, the network has 
    never seen it before. In that way, its quality can be measured against the other methods and is comparable in the
    first place. The other two methods against which \emph{points2mesh} is compared against, should be able to handle
    point cloud data from ModelNet.
    In the following sections, the previously explained metrics are listed and compiled against each other.
    Likewise, \emph{points2mesh}'s reconstructions are listed to comparable methods with the same metrics.
    
\subsection{Distance metric results}
    The first comparison consists of evaluating the distance metric $d_m$ for every configuration
    $\mathcal{C}_i$ and \emph{ball-pivoting algorithm}(~\emph{BPA}~) and \emph{instant field meshes}(~\emph{IFM}~).
    One hundred point clouds of each object category are taken at random from the ModelNet dataset. For each method,
    an approximation mesh $M_{p}$ is calculated, and all distance measure metrics to its ground truth mesh computed.
    The mean value of these metrics is noted in the tables. While the first table denotes reconstruction from 256 
    samples, the following denotes reconstruction from 1024 and subsequently from 7500 points.
    Furthermore, \ref{tab:distance256} also describes the distance metric $d_m$ from the hundred object samples to a not yet deformed sphere\footnote{The sphere object is defined by 
    the implicit surface furmula $x^2+\frac{y^2}{2}+z^2=0$ and created in meshlab with a voxel size of 0.055}.
    This baseline for $d_m$ gives a general understanding of $d_m$. 
    \emph{Deep marching cubes} unfortunately transforms and normalizes their reconstructed mesh in such a way,
    that it cannot be numerically evaluated and thus not compared to \emph{points2mesh} and \emph{IFM} or \emph{BPA}.
    However, a visual comparison is possible and still appealing to conduct, since it is the only other deep learning based
    reconstruction method working (~indirectly\footnote{\emph{DMC} preprocesses the point cloud by transforming it into a voxelized representation, thus not a direct
    tranformation from point cloud data to meshes.}~) on point cloud data, transforming it into a poly-mesh.


    First, as seen in figure \ref{tab:distance256} through \ref{tab:distance7500}, reconstructed meshes already differ depending on the employed
    configuration $\mathcal{C}_i$. For each cell, 100 objects of the \emph{ModelNet} test dataset are used for inference and evaluate
    with $d_m$ against the ground truth. Denoted is the mean value of that distance value. 
    The columns indicate the used configuration or method for the reconstruction. 
    The last column describes the mean value of $d_m$ of a basic spheroidal shape evaluated against 
    all test samples. Comparing this value offers a baseline value, which is an approximate upper bound
    for the quality of the reconstructed meshes. A perfectly reconstructed mesh has a $d_m$ value of $0$.

    The initial goal is to figure out the \emph{best} configuration out of the four proposed ones.
    Reaching that state requires preceding evaluation of $d_m$. Each of them has its merits, but finding the configuration,
    generally capable of the best reconstructions is essential. 
    Taking a look at the mean values of $d_m$ in table \ref{tab:distance1024} for each configuration $\mathcal{C}_i$  shows a small 
    advantage of $\mathcal{C}_1$ and $\mathcal{C}_4$ over the other two configurations. 
    Though small, smaller values are still visible throughout. 
    Exemplary, in figure \ref{fig:1}, the resemblance of the reconstructions are pretty close, though with noticeable differences.
    Taking a closer look at the airplane wings in the first row, $\mathcal{C}_1$, $\mathcal{C}_2$ and $\mathcal{C}_4$ reconstruct
    the small winglets correctly at the end of each wing blade. With configuration $\mathcal{C}_3$, they are not noticeable anymore.
    The advantage can also be seen in the reconstruction of chairs with 1024 samples.  Configuration $\mathcal{C}_1$, $\mathcal{C}_2$
    and $\mathcal{C}_4$ do not have the problem with edges in empty space like configuration $\mathcal{C}_3$, as seen in figure \ref{fig:2}.
    The general low values of $\mathcal{C}_2$, with a higher amount of vertices for the reconstructed object can also be confirmed visually,
    if compared to \ref{fig:1c}, \ref{fig:2c} and \ref{fig:3c}. Though representing the features the best, the general smoothness factor of the object is 
    way lower than the other configurations. Configurations $\mathcal{C}_1$, $\mathcal{C}_3$ and $\mathcal{C}_4$ seem visually more pleasing caused by the smoothness of the surface, though 
    $\mathcal{C}_2$ is able to represent the objects in a better way, if the smoothness is disregarded.

    Furthermore, comparing the results of $d_m$ between 256, 1024 and 7500 samples of each configuration, a more significant gap 
    between 256 to 1024 is present, than between 1024 and 7500 samples.
    While the mean value of $d_m$ gets lower using 1024 samples rather than 256, this is generally not true anymore for changing 
    from 1024 samples to 7500. Since solving the problem of reconstruction is objectively harder the more points are present. 
    However, solving the same problem with only 256 samples may also lead to higher values for $d_m$, since at some amount of samples, there are not enough present
    anymore to represent the ground truth mesh at an appropriate level.
    Notably, in the case of configuration $\mathcal{C}_1$, its mean value for $d_m$ for all object categories is $0.020$ for 256 samples,
    reducing to $0.017$ at 1024 samples, but increasing again to $0.024$ for 7500 samples. Additionally, configuration $\mathcal{C}_3$
    has worse results at 7500 samples than at 1024 samples as well.

    Confirming these numerical results with the visual evidence in figure \ref{fig:7}, 
    while a big jump in visual quality is evident, changing the sample size from 256 to 1024, 
    the change to 7500 samples not so much. As seen in \ref{fig:7d}, the airplane even loses a 
    feature (~fron wheel~) at the last step and thus could explain for higher or equal mean values 
    of $d_m$.

    Tables \ref{distance256opt}, \ref{distance1024opt} and \ref{distance7500opt} show the same evaluation of $d_m$ for $\mathcal{N}_{opt}$, though only the mean value
    is indicated for each configuration. As seen in the reconstructions in figures \ref{10}, \ref{11} and \ref{12}, the results are generally worse. The reconstructions
    have less smooth surfaces, less distinct features shown, and higher values for $d_m$. A complete evaluation for $d_m$ can be found in the appendix \ref{completeopt}

\begin{center}
    \captionof{table}{Distance metric evaluations with 256 samples per point cloud for $\mathcal{N}_{recon}$. Caluclated with \emph{meshdistance} metric. Only considering configurations $\mathcal{C}_1$ and $\mathcal{C}_2$. Lower values are better.} \label{tab:distance256} 
    \begin{tabular}{| l  | c  | c | c | c | c | c || c |}
        \hline
        %airplane
        Object& $\mathcal{C}_1$ & $\mathcal{C}_2$ & $\mathcal{C}_3$ & $\mathcal{C}_4$ & BPA & $IFM$  & None\\ \hline
        airplane&0.013&0.009&0.011&0.011&0.013&0.153&0.251\\\hline
        bed&0.019&0.014&0.018&0.017&0.016&0.123&0.163\\\hline 
        bottle&0.011&0.008&0.011&0.010&0.007&0.134&0.201\\\hline
        bowl&0.024&0.018&0.018&0.022&0.022&0.133&0.148\\\hline
        car&0.020&0.015&0.019&0.020&0.020&0.199&0.155\\\hline
        chair&0.019&0.013&0.018&0.016&0.027&0.188&0.195\\\hline
        guitar&0.015&0.009&0.015&0.013&0.005&0.325&0.338\\\hline
        toilet&0.043&0.020&0.052&0.027&0.016&0.135&0.138\\\hline\hline
        mean&0.020&0.013&0.021&0.016&0.016&0.200&0.1986\\\hline
    \end{tabular}
\end{center}
\begin{center}
    \captionof{table}{Distance metric evaluations with 1024 samples per point cloud for $\mathcal{N}_{recon}$. Caluclated with \emph{meshdistance} metric. Lower values are better.} \label{tab:distance1024} 
    \begin{tabular}{| l  | c | c | c | c | c | c || c |}
        \hline
        %airplane
        Object& $\mathcal{C}_1$ & $\mathcal{C}_2$ & $\mathcal{C}_3$ & $\mathcal{C}_4$ & BPA & $IFM$ & None\\ \hline
        airplane&0.011&0.006&0.012&0.011&0.089&0.10 & 0.251\\\hline
        bed&0.016&0.010&0.016&0.014&0.008&0.099 & 0.163\\        \hline
        bottle&0.009&0.007&0.011&0.008&0.003&0.10 & 0.201\\        \hline
        bowl&0.022&0.013&0.020&0.019&0.006&0.097& 0.148\\        \hline
        car&0.019&0.012&0.018&0.018&0.009&0.163& 0.155\\        \hline
        chair&0.016&0.009&0.017&0.015&0.008&0.172& 0.195\\        \hline
        guitar&0.013&0.006&0.022&0.012&0.002&0.234 & 0.338\\        \hline
        toilet&0.037&0.013&0.048&0.024&0.006&0.109 & 0.138\\        \hline\hline
        mean&0.017&0.010&0.021&0.015&0.016&0.134 & 0.1986\\        \hline
    \end{tabular}
\end{center}
\begin{center}
    \captionof{table}{Distance metric evaluations with 7500 samples per point cloud. Caluclated with \emph{meshdistance} metric.Only considering configurations $\mathcal{C}_1$ and $\mathcal{C}_2$. Lower values are better.} \label{tab:distance7500} 
    \begin{tabular}{| l  | c | c | c | c | c | c |}
        \hline
        %airplane
        Object& $\mathcal{C}_1$ & $\mathcal{C}_2$ & $\mathcal{C}_3$ &$\mathcal{C}_4$ & BPA & $IFM$  \\ \hline
        airplane&0.015&&0.031&0.011&0.090&0.090\\\hline
        bed&0.021&&0.10&0.015&0.005&0.089\\\hline
        bottle&0.014&&0.115&0.007&0.001&0.096\\\hline
        bowl&0.026&&0.127&0.017&0.003&0.091\\\hline
        car&0.024&&0.111&0.018&0.003&0.153\\\hline
        chair&0.022&&0.064&0.016&0.001&0.126\\\hline
        guitar&0.022&&0.025&0.012&0.001&0.212\\\hline
        toilet&0.047&&0.178&0.031&0.003&0.101\\\hline\hline
        mean&0.024&&0.083&0.016&0.013&0.119\\\hline
    \end{tabular}
\end{center}

\begin{center}
    \captionof{table}{Distance metric evaluations with 256 samples per point cloud for $\mathcal{N}_{opt}$. Caluclated with \emph{meshdistance} metric. Only considering configurations $\mathcal{C}_1$ and $\mathcal{C}_2$. Lower values are better.} \label{tab:distance256opt} 
    \begin{tabular}{| l  | c  | c | c | c || c |}
        \hline
        %airplane
        Object& $\mathcal{C}_1$ & $\mathcal{C}_2$ & BPA & $IFM$  & None\\ \hline
        mean&0.035&0.061&0.016&0.200&0.1986\\        \hline
    \end{tabular}
\end{center}
\begin{center}
    \captionof{table}{Distance metric evaluations with 1024 samples per point cloud for $\mathcal{N}_{opt}$. Caluclated with \emph{meshdistance} metric. Lower values are better.} \label{tab:distance1024opt} 
    \begin{tabular}{| l  | c | c | c | c | c | c || c |}
        \hline
        %airplane
        Object& $\mathcal{C}_1$ & $\mathcal{C}_2$ & $\mathcal{C}_3$ & $\mathcal{C}_4$ & BPA & $IFM$ & None\\ \hline
        mean&0.018&0.024&0.021&0.018&0.016&0.134 & 0.1986\\        \hline
    \end{tabular}
\end{center}
\begin{center}
    \captionof{table}{Distance metric evaluations with 7500 samples per point cloud. Caluclated with \emph{meshdistance} metric.Only considering configurations $\mathcal{C}_1$ and $\mathcal{C}_2$. Lower values are better.} \label{tab:distance7500} 
    \begin{tabular}{| l  | c | c | c | c  | c | c |}
        \hline
        Object& $\mathcal{C}_1$ & $\mathcal{C}_2$ & $\mathcal{C}_3$ &  $\mathcal{C}_4$ & BPA & $IFM$  \\ \hline
        mean&0.018&0.015 & 0.025 & 0.027&0.013&0.119\\\hline
    \end{tabular}
\end{center}
    In contrast, the traditional methods, as in \emph{BPA} and \emph{IFM}, get better results; the more samples are provided, both
    numerically as well as visually.
    \emph{IFM} seems to generally struggle to reconstruct visually pleasing results when the number of samples is below 7500.
    Not only that but its numerical values of $d_m$ are sometimes higher than a basic spheroid shape in contrast to its baseline 
    object (~See table \ref{tab:distance256} and \ref{tab:distance1024},  mean value for \emph{car}~). Generally, the metrics have
    higher values for \emph{IFM}, while \emph{BPA} and $\emph{N}_{recon}$ are much closer to each other.

    Not only are they closer to each other, but rather \emph{BPA} produces lower mean values than any configuration of  $\emph{N}_{recon}$ with the exception of $\mathcal{C}_2$.
    Though implicating better results for any case, visually for 256 samples and 1024 samples, $\mathcal{C}_1$ looks in some cases more
    organic, or closer to the ground truth than \emph{BPA} (~Compare figures \ref{fig:7b} with \ref{fig:7c}, and figures \ref{fig:1c} 
    with \ref{fig:1g}~). However, most reconstructions with either \emph{BPA} or \emph{IFM} have holes in the meshes, thus are not watertight.


    Finally, inspecting reconstructions from \emph{DMC} demonstrates generally bad results. The low amount of samples used, less than 7500,
    seems not to be enough for the deep learning based approach to reconstruct a precise object. At figure \ref{fig:51d}, \emph{DMC} somewhat
    reconstruct the general outline of the object, though producing big holes in the side of the object, while also hollowing out the interior
    of the car. Additionally, reconstructing a guitar, see figure \ref{fig:3h}, a second, independent \emph{blob} of incoherent triangles is 
    produced, unrelated to the ground truth. Lastly, as seen in figure \ref{fig:1i}, reconstructing flat surfaces, like the wings of the airplane
    is not possible for \emph{DMC}, resulting in big chunks of missing surface information. However, reconstructed objects with \emph{DMC} are watertight.

\subsection{Generalization capabilities}
    Making use of reconstruction techniques often requires the method to work in a general setting.
    The previous experiments already showed reconstruction from \emph{points2mesh} of test samples,
    which the network has not yet seen. With the same idea of feeding unknown data to the network,
    the next table of evaluated metrics was created. By extending the generalization concept, 
    point cloud samples of objects outside of the trained object categories are fed into the network.
    For that reason, assuming to work on configuration $\mathcal{C}_1$, it has been trained on four different
    object category collections (~See section \ref{sec:categ}~).
    Every trained instance of $\mathcal{C}_1$, now is fed one hundred point cloud samples 
    of \emph{bathtub} and \emph{person}, evaluating every distance metric. Even though, $\mathcal{C}_1$ has never seen any samples of these object categories.


\begin{center}
    \captionof{table}{Distance metric evaluations with 1024 samples per point cloud. Testing with point clouds from outside of trained object categories} \label{tab:generalize} 
    \begin{center}
        \begin{tabular}{| l  | c | c | c | c || c |}
            \hline
            %airplane
            Object& $\mathcal{C}_1$ & $\mathcal{C}_2$ & $\mathcal{C}_3$ & $\mathcal{C}_4$ & None\\ \hline
            bathtub&0.021&0.015&0.067&0.024 & 0.178\\\hline
            person&0.154&0.137&0.181&0.146 & 0.352\\\hline
        \end{tabular}
    \end{center}
\end{center}

    As seen in table \ref{tab:generalize}, all configurations except for $\mathcal{C}_3$ result in low values for $d_m$ if
    compared to the same metrics of no reconstruction at all. Figures \ref{fig:8} and \ref{fig:9} confirm this contention.
    $\mathcal{C}_2$ keeps the same problem as described above, the general roughness of the surface. $\mathcal{C}_1$ and
    $\mathcal{C}_4$ produce rather similar results, while $\mathcal{C}_3$ seems to struggle the most to construct a proper mesh.

\subsection{Varying basemesh $\mathcal{M}_i$}
    The previously described configurations all are trained on a spheroid as initial mesh $\mathcal{M}_i$. 
    By nature of the neural network and the sphere, objects with \emph{genus} higher than zero cannot be
    reconstructed such that the predicted mesh has the same \emph{genus}. Since it stays constant, 
    changing the initial mesh to a toroidal mesh may lead to possible reconstructions with \emph{genus} one. 
    For the following evaluation, $\mathcal{C}_{1}$ is trained on such an initial toroidal mesh.

\begin{center}
    \captionof{table}{Excerpt of evaluation table. Distance metric evaluations with 1024 samples per point cloud. Initial mesh (~torus mesh~) with a \emph{genus} of one.} \label{tab:basemesh} 
    \begin{center}
        \begin{tabular}{| l  | c | c |}
            \hline
            %airplane
            Object& $\mathcal{C}_1^{torus}$ & $\mathcal{C}_1^{sphere}$ \\ \hline
            airplane&0.019&0.011\\\hline
            toilet&0.032&0.029\\\hline
        \end{tabular}
    \end{center}
\end{center}
    Though resulting in only marginally higher values of $d_m$ (~See table \ref{tab:basemesh}~), starting with a toroidal base mesh does not solve the problem 
    of reconstructing objects with a \emph{genus} higher than 0. The final prediction of $\mathcal{N}_{recon}$, assuming the 
    new base mesh, does not show any signs of a hole anywhere in the mesh, no matter the input point cloud. 

\subsection{Noise robustness}
    Noise often can be found in point cloud data in real-world scenarios, 
    since recorded data is subject to imperfections in the recording equipment.
    Thus, reconstructing objects, even with noisy data, is an interesting subject 
    for $\mathcal{N}_{recon}$. Table \ref{tab:noise} shows similar to the earlier 
    experiments, reconstruction of different objects categories with point cloud data,
    with added jitter noise (~Noise level of 0.01~).
    While $d_m$ metrics are only slightly higher than without noise, the visual evidence 
    shows a noticeable difference. Figure \ref{fig:81} depicts reconstructed airplanes with
    jitter noise in the input point cloud data. The general reconstruction still achieves 
    the approximate contour of the object, however having more problems with edges dangling 
    in empty space.
\begin{center}
    \captionof{table}{Excerpt of evaluation table. Distance metric evaluations with 1024 samples per point cloud. Comparison of configuration $\mathcal{C}_1$ trained with and without noise (\emph{noise level} of 0.01). Inference with noisy data of the same \emph{noise level}.} \label{tab:noise} 
    \begin{center}
        \begin{tabular}{| l  | c | c |}
            \hline
            %airplane
            Object& $\mathcal{C}_1^{noisy}$ & $\mathcal{C}_1^{clean}$ \\ \hline
            airplane&0.090&0.012\\\hline
            bed&0.017&0.011\\\hline
            bottle&0.017&0.016\\\hline
            bowl&0.029&0.025\\\hline
            car&0.023&0.021\\\hline
            chair&0.018&0.018\\\hline
            guitar&0.015&0.014\\\hline
            toilet&0.041&0.016\\\hline\hline
            mean&0.031&0.018\\\hline
        \end{tabular}
    \end{center}
\end{center}
\subsection{Evaluation time}
\todo{hypericum stats and graphics card used}
For real-time applications like guiding autonomous cars through traffic without collision,
 reconstruction should happen virtually instantaneously for it to be advantageous. 
 The computation time may be critical, and thus has been evaluated, as seen in table
  \ref{tab:time}. 
\emph{IFM} split into three configurations. $\emph{IFM}^{500}$ takes 256 samples as input 
and produces a mesh with 500 vertices (~which is approximately the most vertices for the best
 result at 256 samples~). Similar, $\emph{IFM}^{2400}$ and $\emph{IFM}^{10100}$ are defined 
 for 1024 and 7500 samples.
$\hat{t}$ is a mean value of the measurement for 100 reconstructed objects. 
Generally, \emph{IFM} has the fastest computation times at the same amount of samples when compared to $\mathcal{C}_i$ and \emph{BPA}.
\begin{center}
    \captionof{table}{Inference time for configurations $\mathcal{C}$ and computation time for other techniques} \label{tab:time} 
        \begin{tabular}{| c | c |}
            \hline
            %airplane
            method & $\hat{t}$ in seconds \\\hline
            $\mathcal{C}_{1,3}$ & 0.061 \\\hline
             $\mathcal{C}_2$ & 0.15 \\\hline
             $\mathcal{C}_4$ & 0.043 \\\hline
             $IFM^{500}$ & 0.001 \\\hline
             $IFM^{2400}$ & 0.005 \\\hline
             $IFM^{10100}$ & 0.167 \\\hline
             $BPA$ & 0.062 \\\hline
        \end{tabular}
\end{center}



\chapter{Discussion}
\label{sec:dicussion}

ergebnisse objektiv durchgehen
Ergebnisse subjektiver betrachten
zu anderen paper vergleichen
torus schlecht
gibt es einen vorteil gegenueber den andern techniken, die im vergleich aufgetaucht sind?

genus problem

was anders machen
wie erweitern

konzept an sich kritisieren


\appendix
\chapter{Appendix}
\label{sec:appendix}

\subsection*{Reconstructions}
\label{optrecon}

\subsection*{Complete evaluations}
\label{completeopt}

\begin{center}
    \captionof{table}{Distance metric evaluations with 256 samples per point cloud for $\mathcal{N}_{opt}$. Caluclated with \emph{meshdistance} metric. Only considering configurations $\mathcal{C}_1$ and $\mathcal{C}_2$. Lower values are better.} \label{tab:distance256opt} 
    \begin{tabular}{| l  | c  | c | c | c || c |}
        \hline
        %airplane
        Object& $\mathcal{C}_1$ & $\mathcal{C}_2$ & BPA & $IFM$  & None\\ \hline
        airplane&0.013&0.12&0.013&0.153&0.251\\\hline
        bed&0.028&0.049&0.016&0.123&0.163\\\hline 
        bottle&0.016&0.014&0.007&0.134&0.201\\\hline
        bowl&0.044&0.089&0.022&0.133&0.148\\\hline
        car&0.035&0.058&0.020&0.199&0.155\\\hline
        chair&0.038&0.045&0.027&0.188&0.195\\\hline
        guitar&0.027&0.039&0.005&0.325&0.338\\\hline
        toilet&0.082&0.075&0.016&0.135&0.138\\\hline\hline
        mean&0.035&0.061&0.016&0.200&0.1986\\        \hline
    \end{tabular}
\end{center}
\begin{center}
    \captionof{table}{Distance metric evaluations with 1024 samples per point cloud for $\mathcal{N}_{opt}$. Caluclated with \emph{meshdistance} metric. Lower values are better.} \label{tab:distance1024opt} 
    \begin{tabular}{| l  | c | c | c | c | c | c || c |}
        \hline
        %airplane
        Object& $\mathcal{C}_1$ & $\mathcal{C}_2$ & $\mathcal{C}_3$ & $\mathcal{C}_4$ & BPA & $IFM$ & None\\ \hline
        airplane&0.011&0.091&0.014&0.011&0.089&0.10 & 0.251\\\hline
        bed&0.016&0.012&0.018&0.016&0.008&0.099 & 0.163\\        \hline
        bottle&0.010&0.007&0.010&0.008&0.003&0.10 & 0.201\\        \hline
        bowl&0.025&0.03&0.029&0.022&0.006&0.097& 0.148\\        \hline
        car&0.021&0.017&0.024&0.021&0.009&0.163& 0.155\\        \hline
        chair&0.018&0.012&0.020&0.017&0.008&0.172& 0.195\\        \hline
        guitar&0.014&0.01&0.022&0.017&0.002&0.234 & 0.338\\        \hline
        toilet&0.029&0.016&0.028&0.025&0.006&0.109 & 0.138\\        \hline\hline
        mean&0.018&0.024&0.021&0.018&0.016&0.134 & 0.1986\\        \hline
    \end{tabular}
\end{center}
\begin{center}
    \captionof{table}{Distance metric evaluations with 7500 samples per point cloud. Caluclated with \emph{meshdistance} metric.Only considering configurations $\mathcal{C}_1$ and $\mathcal{C}_2$. Lower values are better.} \label{tab:distance7500} 
    \begin{tabular}{| l  | c | c | c | c | c | c | c |}
        \hline
        %airplane
        Object& $\mathcal{C}_1$ & $\mathcal{C}_4$ & BPA & $IFM$  \\ \hline
        airplane&0.015&0.092&0.090&0.090\\\hline
        bed&0.016&0.015&0.005&0.089\\\hline
        bottle&0.011&0.009&0.001&0.096\\\hline
        bowl&0.021&0.021&0.003&0.091\\\hline
        car&0.019&0.019&0.003&0.153\\\hline
        chair&0.018&0.016&0.001&0.126\\\hline
        guitar&0.017&0.013&0.001&0.212\\\hline
        toilet&0.033&0.028&0.003&0.101\\\hline\hline
        mean&0.018&0.027&0.013&0.119\\\hline
    \end{tabular}
\end{center}


\bibliography{bibliography}

\end{document}

